% !PRINTED 271
\weberpage{271}
%<<CONT
zwei conjugirte Paare, und daher hat $BB_1$ auch die Form (13) und steht in derselben Beziehung zu $AA_1$, wie $B$ zu $A$ und $B_1$ zu $A_1$.

\gesperrt{Daraus geht hervor, dass die Gesammtheit der Substitutionen $B$ eine mit $\Lambda_p$ isomorphe Gruppe bildet, und diese Gruppe wollen wir mit $\Gamma_p$ bezeichnen.}

Setzen wir $\beta$ an Stelle von $\nu\beta$ und bezeichnen mit $\alpha', \beta'$ die zu $\alpha, \beta$ conjugirten Grössen, so können wir die Substitutionen der Gruppe $\Gamma_p$ auch in der einfacheren Form annehmen:
\weq{14}{B = \begin{pmatrix}\alpha, & \beta\\ -\beta', & \alpha'\end{pmatrix}, \quad \alpha\alpha' + \beta\beta' = 1.}
Um von einer dieser Substitutionen $B$ zu der entsprechenden reellen Substitution $A$ überzugehen, leitet man zunächst aus (14) die entsprechende Substitution $A$ her:
\weq{15}{A = \begin{pmatrix}\alpha, & \nu^{-1}\beta\\ -\nu\beta', & \alpha'\end{pmatrix},}
und bildet daraus nach (3) und (7):
\weq{16}{A = RAR^{-1}.}
Trotz des Isomorphismus lässt sich die Gruppe $\Gamma_p$ nicht in $L_p$ oder $\Lambda_p$ transformiren, wenigstens nicht durch Substitutionen, deren Zahlen dem Körper $\mathfrak{E}_{2,p}$ angehören. Man müsste, um die Transformation auszuführen, in einen höheren Körper gehen, in dem \gesperrt{alle} Zahlen von $\mathfrak{E}_{2,p}$ Quadrate sind.

\begin{center}
\S.\,70.\\
\gesperrt{Divisoren der Gruppe $L_p$, deren Grad durch $p$}\\
\gesperrt{theilbar ist.}
\end{center}

Es ist ein Problem von grösstem Interesse, alle Divisoren der Gruppe $L_p$ zu bestimmen. Von der Lösung dieses Problems hängt es ab, in welcher Weise man die Gruppe $L_p$ durch Permutationen von Ziffern darstellen kann, bei welchen algebraischen Gleichungen also diese Gruppen auftreten können (\S.\,28, 2.).

Die cyklischen Gruppen, die in $L_p$ enthalten sind, haben wir in den vorangehenden Paragraphen schon betrachtet. Wir haben gesehen, dass der Grad einer cyklischen Gruppe entweder gleich $p$ oder ein Theiler von $1/2\,(p-1)$ oder von $1/2\,(p+1)$
%CONT>>
